\section{Polarized cofactors and propagation}
\label{uni-sec:cofactor}

Uniform ellipticity is sufficient for propagation of an extremum, but it is not necessary. The range of the cofactor may already connect enough points to force constancy. For $C^2$ potentials, the quadratic form $Q$ in \eqref{uni-eq:Q} is continuous. Its polarization defines $Q^\sharp:T^*X\to TX$ by $\eta(Q^\sharp\xi)=Q(\xi,\eta)$. A reachable set from a point consists of endpoints of absolutely continuous curves starting there with velocity in $\operatorname{im}Q^\sharp$ almost everywhere. A control orbit consists of endpoints of finite concatenations of flows of locally Lipschitz vector fields taking values in that range.

A positive weight satisfying $\ddc(hT)=0$ gives an energy identity controlling $\mathrm{d}w$ everywhere. Without such a weight, a degenerate maximum principle gives propagation only from the extremal sets. This distinction accounts for the different density assumptions on reachable sets and orbits in the theorem.

\begin{theorem}\label{uni-thm:cofactor-intro}
Let $u,v\in C^2(X)$, $A,B\geq0$, and $A^n=B^n$.
\begin{enumerate}
\item If there is a continuous $h>0$ such that
$\ddc(hT)=0$ in distributions, then
\begin{equation}
  \mathrm{i}\,\partial w\wedge\bar\partial w\wedge T\equiv0.
  \label{uni-eq:energyzero}
\end{equation}
It follows that $w$ is constant if the strict positivity locus of the
cofactor is dense, or if one cofactor control-reachable set is dense.
\item Without the weight hypothesis, the global maximum and minimum sets of
$w$ are invariant under every locally Lipschitz vector field taking values
in the range of the cofactor symbol.  If every such control orbit is dense,
then $w$ is constant.  For smooth data it is enough that smooth range
fields, for example the cofactor columns, and their iterated Lie brackets
span $TX$ at every point.
\end{enumerate}
\end{theorem}

\subsection{Weighted energy}

Testing the weighted cofactor against $w^2$ produces a nonnegative gradient integral. Only distributional $\ddc$-closedness is needed: for a continuous $h>0$, the condition is
\begin{equation}\label{uni-eq:weightedclosed}
 \int_X\ddc\varphi\wedge hT=0\qquad(\varphi\in C^\infty(X,\mathbb R)).
\end{equation}

\begin{proof}[of Theorem~\ref{uni-thm:cofactor-intro}(1)]
The product rule gives
\begin{equation}\label{uni-eq:w2}
 \ddc(w^2)=2w\ddc w+2\mathrm{d}w\wedge \dc w
          =2w\ddc w+\frac{\mathrm{i}}{\pi}\partial w\wedge\bar\partial w.
\end{equation}
Approximate $w^2$ in $C^2$ by smooth functions in
\eqref{uni-eq:weightedclosed}.  Since $hT$ is continuous and
$\ddc w\wedge T=0$, this gives
\begin{equation}\label{uni-eq:energy}
 0=\frac1\pi\int_X h\,\mathrm{i}\partial w\wedge\bar\partial w\wedge T.
\end{equation}
The continuous integrand is nonnegative, so it vanishes pointwise.

For the nonnegative quadratic form $Q$,
$Q(\mathrm{d}w)=0$ implies $Q(\mathrm{d}w,\eta)=0$ for every $\eta$.  Therefore
\begin{equation}\label{uni-eq:annihilate}
 \mathrm{d}w_x\bigl(Q_x^\sharp\eta\bigr)=0\qquad(x\in X,\ \eta\in T_x^*X).
\end{equation}
The chain rule and \eqref{uni-eq:annihilate} make $w$ constant along every absolutely continuous curve whose velocity lies in $\operatorname{im}Q^\sharp$ almost everywhere.  One dense reachable set suffices by continuity.  If
instead $Q$ is positive definite on a dense set, \eqref{uni-eq:energy}
gives $\mathrm{d}w=0$ there and hence everywhere.
\end{proof}

The integration uses only the test function $w^2$. For general tests, integration by parts introduces derivatives of $hT$; $\ddc(hT)=0$ alone does not make $L\varphi=(\ddc\varphi\wedge T)/\mathrm{d}V$ a symmetric divergence operator.

\subsection{The real symbol and exterior-ball propagation}

To propagate an extremum without the weight, we write the linearized operator in real coordinates. Its symbol is continuous and semipositive. An exterior tangent ball to the maximum set can have its normal only in the kernel of that symbol; this forces the maximum set to be invariant under the range fields.

In a holomorphic chart $z_j=x_j+\mathrm{i} y_j$, let $C_{\rm cf}$ be the
Hermitian matrix characterized by
\begin{equation}\label{uni-eq:Ccf}
 \frac{\bigl(\mathrm{i}\sum H_{j\bar k}\mathrm{d}z_j\wedge \mathrm{d}\bar z_k\bigr)\wedge T}{\mathrm{d}V}
     =\operatorname{Tr}(C_{\rm cf}H)
\end{equation}
for every Hermitian matrix $H$.  Positivity of mixed wedges makes
$C_{\rm cf}$ positive semidefinite.  Write
$C_{\rm cf}=P_{\rm cf}+\mathrm{i} Q_{\rm cf}$, where $P_{\rm cf}$ is real
symmetric and $Q_{\rm cf}$ is real skew-symmetric.  In the real coordinate
order $(x_1,\ldots,x_n,y_1,\ldots,y_n)$,
\begin{equation}\label{uni-eq:real-symbol}
 L\varphi:=\frac{\ddc\varphi\wedge T}{\mathrm{d}V}
       =\operatorname{Tr}(aD^2\varphi),\qquad
 a=\frac1{8\pi}
 \begin{pmatrix}P_{\rm cf}&Q_{\rm cf}\\-Q_{\rm cf}&P_{\rm cf}\end{pmatrix}.
\end{equation}
\Needspace{8\baselineskip}
Indeed
$\varphi_{j\bar k}=(\varphi_{x_jx_k}+\varphi_{y_jy_k}
+\mathrm{i}\varphi_{x_jy_k}-\mathrm{i}\varphi_{y_jx_k})/4$; inserting this expression
in \eqref{uni-eq:Ccf} proves \eqref{uni-eq:real-symbol}.  For real coordinate vectors $\xi,\eta\in\mathbb R^n$, moreover,
\[
 (\xi,\eta)^Ta(\xi,\eta)
    =\frac1{8\pi}(\xi-\mathrm{i}\eta)^*C_{\rm cf}(\xi-\mathrm{i}\eta)\geq0.
\]
Thus $a$ is a continuous principal symbol, with no first-order term in
holomorphic coordinates, and its range is that of $Q^\sharp$.

An exterior-ball argument propagates the maximum for this operator.  Let $W\in C^2$
satisfy $LW\geq0$ in a chart, suppose $W\leq M$, and put
$K=\{W=M\}$.  If $z\in K$ is a nearest point to $y\notin K$, then
\begin{equation}\label{uni-eq:normalzero}
 (y-z)^Ta(z)(y-z)=0.
\end{equation}
To prove this, suppose the left side is positive, set
$r_0=|y-z|$, and shrink the tangent ball toward $z$:
$y'=z+\theta(y-z)$, $R=\theta r_0$, with $0<\theta<1$ small.
The closed ball $\overline{B(y',R)}$ lies in the chart and meets $K$
only at $z$.  This follows from the triangle inequality; equality in
$|k-y|\leq|k-y'|+|y'-y|\leq r_0$, for $k\in K$, forces $k=z$.
For large $\gamma$, the function
\[
 \psi(x)=\mathrm{e}^{-\gamma R^2}-\mathrm{e}^{-\gamma|x-y'|^2}
\]
satisfies
\[
 L\psi=2\gamma \mathrm{e}^{-\gamma|x-y'|^2}
   \bigl(\operatorname{tr}a-2\gamma(x-y')^Ta(x)(x-y')\bigr)<0
\]
in a neighborhood of $z$.  Choose a sufficiently small $r<R$, and
consider the lens $D=B(z,r)\cap B(y',R)$.  On its boundary lying on the
outer tangent sphere, $\psi=0$ and $W-M\leq0$.  On the remaining
compact boundary portion, $W-M<0$, since the tangent ball meets $K$
only at $z$.  Hence, for sufficiently small $\varepsilon>0$,
$W-M\leq\varepsilon\psi$ on $\partial D$.  A positive interior
maximum of $W-M-\varepsilon\psi$ would have nonpositive Hessian but
strictly positive image under $L$, which is impossible.  The same
inequality therefore holds in $D$.  Outside the tangent ball
$\psi\geq0$, so it also holds in a full neighborhood of $z$.
Second-order contact at $z$ now yields
\[
 0\leq LW(z)\leq\varepsilon L\psi(z)<0,
\]
proving \eqref{uni-eq:normalzero}.

Since $a\geq0$, its range is the orthogonal complement of its kernel.
Equation \eqref{uni-eq:normalzero} consequently implies
$(y-z)\cdot Z(z)=0$ for every range field $Z$.  If $Z$ is locally
Lipschitz and $\gamma$ is one of its flow curves starting in $K$, let
$d(t)$ denote its squared distance to $K$ in a sufficiently small
compact chart neighborhood.  Shortening the time interval ensures that
nearest points $z_t$ lie in the chart interior.  Write $D^+d(t)=\limsup_{h\downarrow0}(d(t+h)-d(t))/h$ for the upper right derivative. Holding such a nearest
point fixed when differentiating the distance from above gives
\begin{align*}
 D^+d(t)&\leq2(\gamma(t)-z_t)\cdot Z(\gamma(t))\\
 &=2(\gamma(t)-z_t)\cdot\bigl(Z(\gamma(t))-Z(z_t)\bigr)
 \leq2\operatorname{Lip}(Z)d(t).
\end{align*}
As $d(0)=0$, Gronwall's inequality gives $d(t)=0$.  The same argument
for $-Z$, followed by continuation, proves invariance under both signs
of time and under finite concatenations of such flows.

\Needspace{7\baselineskip}
\begin{proof}[of Theorem~\ref{uni-thm:cofactor-intro}(2)]
Apply the preceding argument to $W=w$ and to $W=-w$ in each chart.
Both are solutions of $LW=0$.  The maximum set of $w$ is nonempty,
closed, and contains every orbit starting in it.  If every orbit is dense,
this maximum set is all of $X$, proving constancy.  The argument for the
minimum set is identical.

For completeness, let a smooth family of range fields be bracket
generating.  At $p$, let $\mathcal D_p$ be the span of the values of
all original fields pushed forward by finite compositions of their local
flows and inverse flows.  This distribution contains the original fields
and is invariant under their flows.  Differentiation of pushed-forward
fields along a flow shows that it contains their nested Lie brackets;
finite-dimensional subspaces are closed, so the derivative remains in
$\mathcal D_p$.  Induction and the Jacobi identity give every iterated
bracket.  The assumed rank condition therefore gives
$\mathcal D_p=T_pX$.

Choose $2n$ such pushed-forward fields independent at $p$.  The flow
of each is a conjugate of an original flow and hence preserves its orbit.
The map obtained by successively flowing for $2n$ small times has
invertible derivative at zero.  The inverse function theorem puts a
neighborhood of $p$ in its orbit.  All orbits are open, and their
partition of connected $X$ has only one member.  This proves the stated
H\"ormander-type sufficient condition, without a constant-rank hypothesis
or a smooth matrix square root.  In particular, the columns of a smooth
symbol $a$ may be used as the original range fields.
\end{proof}

For $C^{2,1}$ potentials and a locally Lipschitz background, the cofactor columns are locally Lipschitz and the orbit criterion still applies. The bracket criterion uses smooth fields.
